\section{Manual (Plataforma web) - Apartado de gestión de disciplinas}

Al acceder a la página de disciplinas, se puede visualizar un listado con todas las disciplinas activas que posee el club (de forma similar al apartado de instalaciones y reservas).

\subsection{Detalle de disciplina e inscripción de socios}
Al seleccionar una disciplina específica (por ejemplo, vóley), el sistema permite ver el detalle de la misma. Desde esta pantalla es posible realizar la inscripción directa de los socios.

Pasos para la inscripción:
\begin{itemize}
    \item Seleccionar la disciplina deseada.
    \item Ingresar los datos del socio (por ejemplo, socio número 1001) para efectuar la inscripción.
    \item Para verificar la acción, se puede acceder a la página de socios, donde la nueva disciplina aparecerá reflejada en el perfil del socio inscripto.
\end{itemize}

\subsection{Creación de nuevas disciplinas}
Desde la pantalla principal de disciplinas también es posible dar de alta una nueva actividad. Al momento de crearla, el sistema solicitará configurar los siguientes parámetros:

\begin{itemize}
    \item \textbf{Nombre:} La denominación de la actividad (por ejemplo, prueba).
    \item \textbf{Categoría del socio:} A qué grupo está dirigida (por ejemplo, vitalicios).
    \item \textbf{Sede:} El lugar físico donde se desarrollará.
    \item \textbf{Cupo máximo:} Límite de inscriptos permitidos. Se puede dejar sin límite si la actividad no lo requiere.
    \item \textbf{Arancel:} Se debe indicar si la disciplina es arancelada o gratuita.
\end{itemize}

Una vez completados los datos y confirmada la creación, la nueva disciplina aparecerá inmediatamente en el listado general con todas las características configuradas.
